
% Exposé V: back to the volume's parf
\renewcommand{\parf}{\operatorname{parf}}
\section*{\underline{Exposé V}\\[3pt] \underline{\underline{GÉNÉRALITÉS SUR LES @@M0@@-ANNEAUX}}}
\begin{center}
\underline{par P. BERTHELOT}
\end{center}

La théorie des @@M1@@-anneaux est utile chaque fois que l'on utilise des anneaux tels que l'anneau des classes de fibrés vectoriels sur un site annelé, l'anneau des représentations d'un groupe, etc..., pour lesquels les opérations de puissance extérieure définissent des applications @@M2@@. Nous l'utiliserons abondamment dans la suite du Séminaire, il est donc nécessaire d'en développer d'abord le formalisme général. Cet exposé est donc tout à fait indépendant du reste du séminaire et peut servir aux utilisateurs de @@M3@@-anneaux en dehors de la géométrie algébrique.

Le premier paragraphe donne essentiellement la théorie des « suites multiplicatives de polynômes » de Hirzebruch (cf. [2]); le point de vue fonctoriel adopté ici a l'avantage, outre sa généralité, de réduire notablement les calculs. Les paragraphes suivants montreront comment, à tout @@M4@@-anneau augmenté, on peut associer un « anneau de Chern » permettant de développer une théorie des classes de Chern. C'est cette théorie qui nous permettra d'énoncer le théorème de Riemann--Roch sous la forme générale de Exp. VIII.

\section*{1. Polynômes universels}

1.1. Soit @@M5@@ un anneau commutatif à élément unité, fixé dans tout ce paragraphe. On désigne par @@M6@@ la catégorie des @@M7@@-algèbres filtrées, commutatives à élément unité, dont la filtration est décroissante et telle que le groupe des éléments de filtration nulle soit l'algèbre entière. On définit un foncteur sur @@M8@@, à valeurs dans la catégorie des groupes 
